% ===== 7 实验结果与分析 =====
\section{实验结果与分析}
\label{sec:results}

本章按照第\ref{sec:experimental-design}章规定的证据等级组织结果。当前仓库已经具备三张自然堆积骨料图像的完整端到端推理产物、规则模块输出以及合成单元测试，因此可以回答三个问题：处理链能否在真实场景图像上形成稳定输出；数量统计、质量代理、异常剔除等设计会怎样改变结果；软件实现是否与预期数学定义一致。与此同时，当前仓库尚未提交与三张演示样本配对的标准机械筛分原始质量，因此本章\textbf{不将演示样本的预测值解释为筛分准确率}。与机械筛分真值有关的最终误差表和结论将在配对物理验证完成后回填。

\subsection{功能演示集与结果口径}
\label{sec:demo-results-scope}

当前功能演示集包括 \texttt{agg\_001}、\texttt{agg\_005} 和 \texttt{agg\_029} 三张自然堆积场景图像。三张图像均使用已经写入推理产物的尺度 $0.208$~mm/px，并在当前基线参数 $\boldsymbol{\theta}=(1,0,0)$、$\gamma=3$ 下生成实例级几何表、筛分等效粒径、质量代理级配、形貌标签、异常标签和可视化报告。由于 $\boldsymbol{\theta}$ 与 $\gamma$ 尚未通过本赛题配对机械筛分批次校准，表\ref{tab:demo-main-results}中的 $D_p$ 应理解为\emph{当前视觉质量代理模型的预测输出}，而不是已经验证的机械筛分真值。

三张图像分别检测到 1924、968 和 966 个实例；按照当前异常规则保留为正常骨料的实例数分别为 1193、703 和 520。三者形成了明显不同的预测级配形态：\texttt{agg\_001} 的正常实例质量代理主粒级为 25--31.5~mm，\texttt{agg\_005} 同样由 25--31.5~mm 主导但粗粒尾部更强，\texttt{agg\_029} 则由 5--10~mm 主导。由此可见，同一处理链能够把不同视觉场景转换为不同的场景级粒度描述，而不是输出近似固定的分布。

\begin{table}[htbp]
\centering
\caption{三张功能演示样本的当前视觉质量代理输出。$D_p$ 与粒级占比均基于当前未校准参数，不能单独解释为机械筛分准确率。}
\label{tab:demo-main-results}
\small
\renewcommand{\arraystretch}{1.15}
\begin{tabular}{lrrrrll}
\toprule
样本 & 实例数 & 正常实例 & $D_{10}$ / mm & $D_{50}$ / mm & $D_{90}$ / mm & 正常实例主粒级 \\
\midrule
\texttt{agg\_001} & 1924 & 1193 & 11.1 & 23.6 & 30.6 & 25--31.5~mm（34.2\%） \\
\texttt{agg\_005} & 968  & 703  & 15.6 & 27.5 & 38.0 & 25--31.5~mm（29.7\%） \\
\texttt{agg\_029} & 966  & 520  & 6.7  & 14.6 & 30.6 & 5--10~mm（32.7\%） \\
\bottomrule
\end{tabular}
\end{table}

图\ref{fig:demo-output-example}以 \texttt{agg\_001} 为例展示从实例检测结果到场景级粒度分布的两个输出层次。左图用于检查实例边界、漏检和粘连等上游视觉问题，右图用于检查数量型与质量代理型累计分布之间的统计差异。将两者并列的目的，是保持场景级级配结果对上游实例测量可追溯：一旦最终级配出现异常偏粗或偏细，可以回到实例层定位是否由 merge、split、边缘截断或异常目标引起，而不是只观察一个最终 $D_{50}$ 数字。

\begin{figure}[htbp]
\centering
\begin{subfigure}{0.48\textwidth}
  \centering
  \includegraphics[width=\linewidth]{agg_001_IG2_full_set_cp_SAM_pred_grains_re_scaled_detection_overview.png}
  \caption{实例检测与粒径标注概览}
\end{subfigure}\hfill
\begin{subfigure}{0.48\textwidth}
  \centering
  \includegraphics[width=\linewidth]{agg_001_IG2_full_set_cp_SAM_pred_grains_re_scaled_gsd_comparison.png}
  \caption{数量分布与质量代理分布比较}
\end{subfigure}
\caption{\texttt{agg\_001} 的端到端功能演示。该图证明处理链能够从实例结果生成级配输出，但没有机械筛分真值叠加，因此不用于声明预测准确率。}
\label{fig:demo-output-example}
\end{figure}

\subsection{数量统计与质量代理的差异}
\label{sec:number-vs-mass-results}

当前三张样本最稳定、也最重要的行为性观察，是\textbf{数量型粒径分布与质量代理型粒径分布存在显著差异}。表\ref{tab:number-mass-comparison}给出了异常处理前、基于全部检测实例计算的两种分位数。以 $D_{50}$ 为例，\texttt{agg\_001} 从数量口径的 6.0~mm 移动到质量代理口径的 23.5~mm；\texttt{agg\_005} 从 7.6~mm 移动到 27.5~mm；\texttt{agg\_029} 从 5.2~mm 移动到 14.1~mm。三张样本方向一致，说明场景中大量小实例在\emph{颗粒个数}上占优，而少量较大实例在 $d^3$ 权重下对\emph{总质量代理}贡献更大。

\begin{table}[htbp]
\centering
\caption{数量型与质量代理型分位数的行为比较。质量代理采用当前默认 $w=d^3$；该比较用于说明统计口径影响，不用于证明 $\gamma=3$ 的物理准确性。}
\label{tab:number-mass-comparison}
\small
\renewcommand{\arraystretch}{1.15}
\begin{tabular}{lrrrrrr}
\toprule
& \multicolumn{3}{c}{数量型分位数 / mm} & \multicolumn{3}{c}{质量代理分位数 / mm} \\
\cmidrule(lr){2-4}\cmidrule(lr){5-7}
样本 & $D_{10}$ & $D_{50}$ & $D_{90}$ & $D_{10}$ & $D_{50}$ & $D_{90}$ \\
\midrule
\texttt{agg\_001} & 3.1 & 6.0 & 17.3 & 10.6 & 23.5 & 30.6 \\
\texttt{agg\_005} & 3.3 & 7.6 & 24.5 & 15.4 & 27.5 & 38.0 \\
\texttt{agg\_029} & 3.0 & 5.2 & 9.2  & 5.9  & 14.1 & 30.6 \\
\bottomrule
\end{tabular}
\end{table}

这一结果与第\ref{sec:mass-weighting}节的机制一致。数量型经验分布给每个实例相同权重，因此一个 4~mm 颗粒和一个 30~mm 颗粒各贡献一个计数；质量代理在 $\gamma=3$ 时二者权重比约为 $(30/4)^3$，使粗颗粒对累计分布产生更大影响。因此，$D_{50}$ 向粗粒方向移动并非后处理中的偶然数值变化，而是由“从颗粒数量统计切换到质量相关统计”的定义直接导致。

但这一行为性一致并不能推出 $\gamma=3$ 已经最接近机械筛分。$\gamma$ 越大，粗粒权重天然越高，$D_p$ 也会倾向向右移动；只有当该移动使独立测试批次相对于机械筛分真值的 $E_{D_p}$、$\mathrm{MAE}_{\mathrm{sieve}}$ 和逐筛级误差下降时，才能说质量模型得到物理验证。因此，当前结果支持的是“\emph{质量口径不可忽略}”这一设计判断，而“\emph{最佳质量指数是多少}”仍属于第\ref{sec:sieve-calibration}节需要通过配对筛分批次解决的问题。

\subsection{粒级组成与异常剔除的影响}
\label{sec:grading-anomaly-results}

三张演示样本进一步显示，不同场景的主粒级和粗细尾部可以通过同一套筛孔边界进行结构化表达。\texttt{agg\_001} 在正常实例口径下，25--31.5~mm 质量代理占比为 34.2\%，20--25~mm 为 21.1\%；\texttt{agg\_005} 的 25--31.5~mm 与 31.5--40~mm 分别约为 29.7\% 和 25.9\%，因而呈现更明显的粗粒组成；\texttt{agg\_029} 的 5--10~mm 占比为 32.7\%，10--16~mm 为 19.1\%，呈现相对细粒的预测结构。这些差异与表\ref{tab:demo-main-results}中的 $D_{50}$ 排序一致：\texttt{agg\_029}<\texttt{agg\_001}<\texttt{agg\_005}。

当前异常分类结果见表\ref{tab:demo-anomaly-results}。三张样本均未触发 $d>50$~mm 的 oversized 规则，也未触发当前 40--50~mm 且低 solidity 的疑似泥团规则；主要被标记的异常是 $d<5$~mm 的 undersized 实例，其数量占比为 27.4\%--46.2\%。这说明在实例数量口径下，小尺度目标非常常见，来源可能同时包含真实细小颗粒、边界碎片、局部分割残片和尺度接近下限的检测结果，不能仅凭“数量很多”判断其物理质量贡献很大。

\begin{table}[htbp]
\centering
\caption{三张功能演示样本的异常规则输出。该表反映当前规则触发情况，不等价于异常检测 Precision/Recall。}
\label{tab:demo-anomaly-results}
\small
\renewcommand{\arraystretch}{1.15}
\begin{tabular}{lrrrr}
\toprule
样本 & 正常 & $>50$~mm & $<5$~mm & 疑似泥团 \\
\midrule
\texttt{agg\_001} & 62.0\%（1193） & 0 & 38.0\%（731） & 0 \\
\texttt{agg\_005} & 72.6\%（703）  & 0 & 27.4\%（265） & 0 \\
\texttt{agg\_029} & 53.8\%（520）  & 0 & 46.2\%（446） & 0 \\
\bottomrule
\end{tabular}
\end{table}

以 \texttt{agg\_001} 为例，全部实例的质量代理 $D_{10}$ 为 10.57~mm，剔除当前异常规则后变为 11.12~mm；$D_{50}$ 从 23.50~mm 变为 23.56~mm，而 $D_{90}$ 基本不变。其原因是 $<5$~mm 实例虽然占检测数量的 38.0\%，但在 $d^3$ 权重下对总质量代理的贡献有限。因此，当前数据支持“数量很多的小目标不一定主导质量级配”这一机制判断，也说明异常剔除对场景级结果的影响与异常类别和粒径位置有关。

相反，粗尺度误检的风险更高。根据式\eqref{eq:mass-weight}，错误 merge 或 oversized 目标的权重会随 $d^\gamma$ 快速增长，因此极少量粗粒异常即可对高分位数产生不成比例的影响。这也是后续独立测试中必须将 merge 错误、真实超限颗粒和非骨料异物分别统计的原因。当前三张演示图没有触发 $>50$~mm 规则，因而不能用于评价该规则对真实异物的召回率；这一结论必须依赖第\ref{sec:ground-truth}节定义的含已知异常目标测试集。

\subsection{投影形貌输出}
\label{sec:shape-demo-results}

表\ref{tab:demo-shape-results}给出了当前几何阈值在三张演示图上的投影形貌分布。三张图中“普通”类别均占多数，为 55.0\%--64.1\%；针片状候选约为 14.0\%--20.9\%；圆形和棱角状的比例随场景发生变化。该结果至少说明形貌模块能够对实例级几何特征产生非退化的多类别输出，并且不同场景不会得到完全相同的类别比例。

\begin{table}[htbp]
\centering
\caption{三张功能演示样本的二维投影形貌规则输出（按实例数量统计）。}
\label{tab:demo-shape-results}
\small
\begin{tabular}{lrrrr}
\toprule
样本 & 针片状候选 & 圆形 & 棱角状 & 普通 \\
\midrule
\texttt{agg\_001} & 16.3\% & 8.0\%  & 16.4\% & 59.3\% \\
\texttt{agg\_005} & 20.9\% & 6.6\%  & 17.6\% & 55.0\% \\
\texttt{agg\_029} & 14.0\% & 11.0\% & 11.0\% & 64.1\% \\
\bottomrule
\end{tabular}
\end{table}

但是，多类别输出本身不构成分类准确性证据。当前阈值来自显式几何规则，尚缺少独立人工形貌标签与真实三维针片状物理量规，因此这里不能从类别比例推断“某类识别准确”。特别是“针片状候选”只描述俯视投影的细长程度，不能替代厚度维度参与的规范针片状判定。正式结果应在独立标注集上增加混淆矩阵和 Macro-F1，并在需要声称真实针片状识别时引入第三维物理真值。

\subsection{软件一致性与可复现性}
\label{sec:software-validation-results}

除真实图像功能演示外，仓库还为筛分等效粒径、质量权重、形貌、异常和报告模块提供了不依赖 Cellpose 模型的合成测试。当前竞赛层测试覆盖 \texttt{test\_sieve\_equivalent.py}、\texttt{test\_morphology.py}、\texttt{test\_anomaly.py} 和 \texttt{test\_report.py}，仓库测试记录为 27 项通过。测试内容包括等效直径与 $w=d^\gamma$ 计算、加权分位数、合成校准参数恢复、四类形貌阈值、异常优先级以及场景摘要/消融表生成等。

这组测试支持的结论是\textbf{实现行为与代码定义一致}。例如，在已知合成参数生成的数据上，校准函数能够搜索到与设定参数接近的解，说明优化器、损失函数和数据接口在数值层面形成闭环；形貌和异常测试则用于防止阈值修改、列名兼容或优先级逻辑在后续工程迭代中发生回归。它们对于最终现场系统的可复现性是必要条件，但由于输入是合成数据，不能替代真实实例掩膜、机械筛分或人工材质标签的外部有效性验证。

三张演示样本的完整中间产物也已保存在仓库，包括预测掩膜、逐颗粒 CSV、毫米尺度表、形貌/异常标注、级配汇总 JSON 以及多种可视化。因而当前报告中的表\ref{tab:demo-main-results}--\ref{tab:demo-shape-results}可以从提交产物重新读取，而不依赖人工抄录隐藏结果。正式提交前还应将相同的数据留痕原则扩展到机械筛分真值：保存逐筛质量原始记录、批次编号、参数配置和最终误差表，使评委可以从原始称量追溯到报告数字。

\subsection{工程运行时间：现有记录与正式计时要求}
\label{sec:runtime-results}

旧版实验记录给出的单张处理时序约为实例分割 17~s、几何测量 9~s、筛分统计小于 1~s，即算法处理部分约 27~s/张；三张顺序处理约 81~s。该记录表明在当时测试硬件与已有模型缓存条件下，算法计算本身远低于比赛 30~min 时间窗口，可以作为工程可行性的初步证据。

然而，按照第\ref{sec:engineering-evaluation}节的定义，27~s/张不能单独承担“现场总流程小于 30~min”的最终 claim。比赛实际消耗还包括设备摆放、尺度标定、图像采集、图像质量检查、必要的重拍、结果目检和导出。因此，正式现场验证应重新采用统一硬件至少重复若干次完整 SOP，并报告各阶段的 wall-clock time、均值/范围以及总耗时。算法阶段的旧记录可保留为基线，但最终结论应以端到端流程计时为准。

\subsection{标准机械筛分精度：当前证据缺口与回填格式}
\label{sec:physical-accuracy-results}

截至当前版本，仓库中的三张功能演示样本没有配对提交逐筛机械称量数据，因此无法计算式\eqref{eq:metric-d-error}和式\eqref{eq:metric-sieve-mae}定义的独立筛分误差。这个缺口不是排版问题，而是当前报告最关键的实验缺口：没有同批次物理真值，就无法判断当前 $\boldsymbol{\theta}=(1,0,0)$ 与 $\gamma=3$ 是否只产生了视觉上合理的级配，还是确实逼近标准机械筛分。

正式实验完成后，本节应首先报告\textbf{冻结参数后的独立测试结果}，推荐使用表\ref{tab:physical-result-template}的最小格式。若有多个测试批次，应同时给出逐批结果和总体平均/标准差，而不是只挑选表现最好的样本。除三个 $D_p$ 外，还必须报告完整筛孔累计通过率或逐筛级质量占比误差，因为当前校准损失只直接约束 $D_{10}/D_{50}/D_{90}$。

\begin{table}[htbp]
\centering
\caption{独立机械筛分测试结果的回填模板。当前版本暂不填数值，避免将功能演示误写成物理精度。}
\label{tab:physical-result-template}
\small
\renewcommand{\arraystretch}{1.15}
\begin{tabular}{lrrrrrrr}
\toprule
批次 & $D_{10}^{\mathrm{GT}}$ & $\widehat D_{10}$ & $D_{50}^{\mathrm{GT}}$ & $\widehat D_{50}$ & $D_{90}^{\mathrm{GT}}$ & $\widehat D_{90}$ & $\mathrm{MAE}_{\mathrm{sieve}}$ \\
\midrule
Test-01 & -- & -- & -- & -- & -- & -- & -- \\
Test-02 & -- & -- & -- & -- & -- & -- & -- \\
Test-03 & -- & -- & -- & -- & -- & -- & -- \\
\midrule
平均 & -- & -- & -- & -- & -- & -- & -- \\
\bottomrule
\end{tabular}
\end{table}

在同一独立测试集上，还应完成第\ref{sec:ablation-design}节定义的四组消融。真正有解释力的结果不是“$\gamma=3$ 使 $D_{50}$ 变大”，而是“质量模型相对于数量基线使机械筛分误差降低多少”；不是“混合 $b$ 与 $d_{\mathrm{eq}}$ 后数值发生变化”，而是“校准后的混合模型在未见批次上是否稳定优于单一几何量”；也不是“异常剔除后曲线更平滑”，而是“对已确认异常目标的处理是否降低最终级配误差”。只有这样，消融才能从参数敏感性展示升级为方法有效性证据。

\subsection{阶段性结论与下一步实验优先级}
\label{sec:interim-results-conclusion}

综合当前证据，可以形成三项阶段性结论。第一，实例分割、毫米几何、筛分等效粒径、质量代理、形貌与异常模块已经在三张真实自然堆积图像上形成完整的数据闭环，并能够保存可追溯的中间产物。第二，数量型与质量代理型分布在三张样本上均存在显著差异，证明“从实例数量转向质量相关统计”是不能省略的建模环节；但最佳 $\gamma$ 和 $\boldsymbol{\theta}$ 尚需机械筛分真值确定。第三，当前软件测试能够保证核心数学操作和规则逻辑的一致性，但实例识别精度、机械筛分精度、形貌准确率和泥团检测准确率仍分别缺少对应的外部 ground truth。

因此，后续实验的最高优先级不是继续增加更多无真值的场景截图，而是建立若干\textbf{同批图像--机械筛分}配对批次，并严格分离校准与独立测试。第二优先级是为代表性图像补充人工实例掩膜，使最终级配误差可以进一步分解为上游实例误差与下游筛分语义误差。第三优先级才是扩展投影形貌和泥团/异物人工标签。完成这三类真值后，本章现有结构无需重写，只需将表\ref{tab:physical-result-template}及相应消融位置替换为真实独立测试数字，即可从“功能可行报告”升级为“具有物理精度证据的竞赛技术报告”。

在明确当前证据边界后，下一章转向工程实现与现场部署，回答即使算法链本身成立，如何在比赛规定的有限时间内通过稳定成像、尺度校验、自动化执行和结果质检把方法可靠地运行起来。
